\documentclass[11pt]{article}
\usepackage[margin=2.5cm]{geometry}
\usepackage{amsmath,amssymb}
\usepackage{booktabs,tabularx,array}
\usepackage{xcolor}
\usepackage[authoryear,round]{natbib}
\usepackage[colorlinks=true,linkcolor=blue,citecolor=blue,urlcolor=blue]{hyperref}

\newcommand{\E}{\mathbb{E}}
\newcommand{\F}{\mathcal{F}}
\newcommand{\dist}{\operatorname{dist}}
\newcolumntype{Y}{>{\raggedright\arraybackslash}X}

\begin{document}

% ============================================================
\section{Basic Stochastic Convex Optimization: Rates and Lower Bounds for SGD}
% ============================================================

\subsection{Setup and Noise Assumptions}

$f$ is convex and $L$-smooth, $S:=\arg\min f\neq\emptyset$, $x^\star\in S$,
$D:=\|x_0-x^\star\|$, and $\Delta_t:=f(x_t)-f^\star$. SGD \citep{robbins1951} is given by
\[
x_{t+1}=x_t-s_t\bigl(\nabla f(x_t)+e_t\bigr),
\]
where $\E_t[\cdot]:=\E[\cdot\mid\F_t]$ and $e_t$ is the stochastic gradient error.
The literature uses different noise assumptions. We refer to them below by the labels in Table~\ref{tab:noise}.

\begin{table}[htbp]
\centering
\small
\renewcommand{\arraystretch}{1.3}
\caption{Noise assumptions ($g_k=\nabla f(\cdot)+e_k$ is a stochastic gradient)}
\label{tab:noise}
\begin{tabularx}{\textwidth}{@{}p{1.0cm}p{2.6cm}Yp{3.2cm}@{}}
\toprule
Label & Name & Mathematical form & References \\
\midrule
(N1) & Conditional unbiasedness & $\E[e_k\mid\F_k]=0$ & All works \\
(N2) & Bounded variance & $\E[\|e_k\|^2\mid\F_k]\le\sigma^2$ & \citet{lan2012}; JOTA setting \\
(N3) & Bounded gradients & $\|\partial f(x)\|\le G$ or $\|g(x,\xi)\|\le G$ & \citet{nemirovski2009} (smooth) \\
(N4) & ABC condition & $\E\|g_k\|^2\le A\,(f(x_k)-f^\star)+B\,\|\nabla f(x_k)\|^2+C$ & \citet{liu2022} \\
(N5) & Relative noise & $\E[\|e_k\|^2\mid\F_k]\le\nu^2\|\nabla f(y_k)\|^2$ & \citet{vaswani2019fast,gupta2024nesterov,hermant2025gradient} \\
(N6) & Decaying noise & $\sigma_k^2:=\E[\|e_k\|^2\mid\F_k]$, $(s_kt_k\sigma_k)\in\ell^2$ or $\ell^1$ & \citet{attouch2026ravine} \\
\bottomrule
\end{tabularx}
\end{table}

(N2) and (N3) are special cases of (N4). We state (N5) at the gradient evaluation point $y_k$, since stochastic Nesterov evaluates the gradient at the extrapolated point.

\subsection{Basic Stochastic Convex Optimization: Rates and Lower Bounds for SGD}

Under (N1), (N2) and $s_t\le 1/(2L)$, expand $\|x_{t+1}-x^\star\|^2$ and take conditional expectations (see Appendix~\ref{app:sgd}):
\begin{equation}
\E_t\|x_{t+1}-x^\star\|^2+s_t\Delta_t\le\|x_t-x^\star\|^2+s_t^2\sigma^2 .
\label{eq:onestep}
\end{equation}
\textbf{Interpretation:} the distance decreases by at least $s_t\Delta_t$, at the cost of the noise term $s_t^2\sigma^2$.

Let $S_k:=\sum_{t<k}s_t$ and $\widetilde x_k:=S_k^{-1}\sum_{t<k}s_tx_t$. Taking total expectations in \eqref{eq:onestep}, summing from $0$ to $k-1$, and applying Jensen’s inequality gives:
\begin{equation}
\boxed{\;
\E[f(\widetilde x_k)-f^\star]
\le
\underbrace{\frac{D^2}{S_k}}_{\text{bias}}
+\underbrace{\frac{\sigma^2\sum_{t<k}s_t^2}{S_k}}_{\text{variance}} .
\;}
\label{eq:biasvar}
\end{equation}
\textbf{Interpretation:} the bias term decreases as the step size grows, the variance term increases with the step size, and the rate comes from balancing the two.

\subsubsection{Rates from \texorpdfstring{\eqref{eq:biasvar}}{(2)}}

\begin{center}
\small
\begin{tabularx}{\linewidth}{@{}Yl Y@{}}
\toprule
Step size & Expected rate (weighted average) & Description \\
\midrule
constant $s\asymp\min\{1/L,\,D/(\sigma\sqrt T)\}$, with known $T$ & $O\bigl(LD^2/T+\sigma D/\sqrt T\bigr)$ & balancing $D^2/(sT)$ and $s\sigma^2$ \citep{nemirovski2009} \\
$s_t\asymp t^{-p}$, $p<1/2$ & $O(k^{-p})$ & variance \\
$p=1/2$ & $O(\log k/\sqrt k)$ & best choice when $T$ is unknown \\
$p=1/2+\varepsilon$, $\varepsilon\in(0,1/2)$ & $O(k^{-1/2+\varepsilon})$ & $\sum_t s_t^2<\infty$; arbitrarily close to the critical $k^{-1/2}$ rate \\
\bottomrule
\end{tabularx}
\end{center}
The last three rows follow by directly substituting $s_t\asymp t^{-p}$ into \eqref{eq:biasvar}.

\subsubsection{Lower Bound Including the Bias Term}

Under the smoothness assumption and (N1), (N2), the minimax lower bound in the stochastic first-order oracle model is
\[
\Omega\Bigl(\frac{LD^2}{T^2}+\frac{\sigma D}{\sqrt T}\Bigr).
\]
The lower bound for the bias term follows from deterministic first-order methods \citep[Theorem~2.1.7]{nesterov2004};
the lower bound for the noise term follows from information-theoretic arguments \citep{nemirovski1983,agarwal2012}.
With the AC-SA method, both terms are attained simultaneously \citep{lan2012}.
Thus, SGD is already optimal in the noise term; acceleration only improves the bias term from $LD^2/T$ to $LD^2/T^2$.
As $k\to\infty$, the asymptotic order is determined by $\sigma D/\sqrt k$: \textbf{there is no asymptotic acceleration under persistent noise (N2).}

\subsection{Some Remarks on Almost-sure Analysis}

The basic sufficient conditions are
\[
\sum_t s_t=\infty,
\qquad
\sum_t s_t^2<\infty.
\]
Under these conditions, the noise term in \eqref{eq:onestep} is summable, and the Robbins--Siegmund lemma \citep{robbins1971} gives
\[
\sum_t s_t\Delta_t<\infty \qquad \text{a.s.}
\]
For polynomial stepsizes $s_t\asymp t^{-p}$, square summability requires $p>1/2$. Writing $p=1/2+\varepsilon$ with arbitrarily small $\varepsilon>0$, the expectation bound above becomes $O(k^{-1/2+\varepsilon})$. Thus, within this stepsize-based analysis, the rate can approach but does not attain the critical $k^{-1/2}$ rate; a corresponding almost-sure rate requires an additional argument beyond the summability conclusion above.

% ============================================================
\section{Stochastic Nesterov: The Same Structure and the Same Limit}
% ============================================================

Stochastic Nesterov uses
\[
y_k=x_k+\alpha_k(x_k-x_{k-1}),\qquad
x_{k+1}=y_k-s_k(\nabla f(y_k)+e_k).
\]
The momentum sequence is encoded by
\[
t_k:=1+\sum_{i\ge k}\prod_{j=k}^{i}\alpha_j,
\qquad\text{equivalently}\qquad
\alpha_k=\frac{t_k-1}{t_{k+1}},
\]
and we set $z_k:=x_{k-1}+t_k(x_k-x_{k-1})$. JOTA \citep[Theorem~3.1]{attouch2026ravine} assumes that $f$ is convex, $\nabla f$ is $L$-Lipschitz, $S\neq\emptyset$, $(s_k)\subset(0,1/L]$ is nonincreasing, $x_0=x_1$, $\F_k:=\sigma(x_i,e_{i-1}:i\le k)$ and $\sigma_k^2:=\E[\|e_k\|^2\mid\F_k]$. The conditions on the momentum and the noise are:
\begin{description}
\item[$(K_0)$] $\sum_{i\ge k}\prod_{j=k}^{i}\alpha_j<\infty$ for all $k\ge1$, i.e.\ $t_k$ is finite.
\item[$(K_1)$] $t_{k+1}^2-t_k^2\le t_{k+1}$ for all $k\ge1$. This bounds the growth of $t_k$ (it implies $t_{k+1}\le 2t_k$) and makes the function-value term in the one-step inequality below nonpositive. For $\alpha_k=1-\alpha/k$ with $\alpha>1$, one has $t_{k+1}=k/(\alpha-1)$, and $(K_1)$ holds if and only if $\alpha\ge3$.
\item[$(K_1^+)$] There is $m\in[0,1)$ such that $t_{k+1}^2-t_k^2\le m\,t_{k+1}$ for all $k\ge1$. This leaves a strictly negative margin in the one-step inequality, which gives summability and $o$ rates. For $\alpha_k=1-\alpha/k$, $(K_1^+)$ holds if and only if $\alpha>3$ (with $m=2/(\alpha-1)$).
\item[$(K_2)$] $\E[e_k\mid\F_k]=0$ and $\sum_k s_k^2t_k^2\sigma_k^2<\infty$ a.s., i.e.\ $(s_kt_k\sigma_k)\in\ell^2_+(\F)$. This makes the noise term in the one-step inequality summable.
\item[$(K_2^+)$] $\E[e_k\mid\F_k]=0$ and $\sum_k s_kt_k\sigma_k<\infty$ a.s., i.e.\ $(s_kt_k\sigma_k)\in\ell^1_+(\F)$. This is used for iterate convergence and for the transfer to the Ravine sequence $y_k$.
\end{description}

The Lyapunov energy used in JOTA is
\[
V_k=s_k\,t_k^2(f(x_k)-f^\star)+\tfrac12\dist(z_k,S)^2.
\]
Using the unbiasedness in $(K_2)$, one step gives \citep[eq.~(18)]{attouch2026ravine}
\[
\E[V_{k+1}\mid\F_k]\le V_k-s_k\bigl(t_k^2-t_{k+1}^2+t_{k+1}\bigr)(f(x_k)-f^\star)+s_k^2t_{k+1}^2\sigma_k^2 .
\]
Under $(K_1)$, the middle term is nonpositive and $t_{k+1}\le2t_k$, so
\[
\E[V_{k+1}\mid\F_k]\le V_k+4s_k^2t_k^2\sigma_k^2 .
\]
Under $(K_1^+)$, the middle term is at most $-(1-m)s_kt_{k+1}(f(x_k)-f^\star)$ and is kept. Under the summability in $(K_2)$, the Robbins--Siegmund lemma gives that $V_k$ converges a.s., hence $f(x_k)-f^\star=O(1/(s_kt_k^2))$ a.s. Thus the noise enters only through the weighted square term. Taking expectations and summing yields \citep[Theorem~3.1(i)]{attouch2026ravine}
\[
\E[f(x_k)-f^\star]\le
\underbrace{\frac{\E[V_1]}{s_kt_k^2}}_{\text{bias}}+
\underbrace{\frac{4\sum_{i\le k}s_i^2t_i^2\,\E[\sigma_i^2]}{s_kt_k^2}}_{\text{variance}},
\qquad
\E[V_1]=s_1t_1^2\,\E[f(x_0)-f^\star]+\tfrac12\E\bigl[\dist(x_0,S)^2\bigr].
\]
This is the same bias--variance structure as \eqref{eq:biasvar}, with $s_kt_k^2$ replacing $\sum_{t<k}s_t$. Acceleration enlarges the denominator, but it also amplifies the noise by $t_i^2$.

For $t_k\asymp k$, the two principal regimes are:
\begin{center}
\small
\begin{tabularx}{\linewidth}{@{}lYllY@{}}
\toprule
Noise regime & Parameters & Bias & Variance & Consequence \\
\midrule
Decaying noise (N6) & $s$ constant, $\sum_k k^2\sigma_k^2<\infty$ & $1/k^2$ & $1/k^2$ & acceleration is preserved; with $(K_1^+)$, a.s.\ $o(1/k^2)$ \\
Persistent noise (N2) & $\sigma_k\equiv\sigma$, $s_k=c k^{-\delta}$ & $k^{\delta-2}$ & $k^{1-\delta}$ & balance at $\delta=3/2$, giving expected $O(\log k/\sqrt{k})$ \\
\bottomrule
\end{tabularx}
\end{center}

Under persistent noise, almost-sure analysis requires
$\sum_k k^{2-2\delta}<\infty$, hence $\delta>3/2$. The resulting rates are slightly slower than $k^{-1/2}$ and can approach it arbitrarily closely. Therefore stochastic Nesterov and SGD have the same noise-dominated order under persistent noise; acceleration survives when the noise decays.

% ============================================================
\section{Research Landscape of Stochastic Acceleration}
% ============================================================

Results for stochastic acceleration must be compared under the same noise model, evaluation point, convergence notion, and total sample budget. Whether stochastic Nesterov accelerates depends on the noise model; Table~\ref{tab:lines} summarizes the three cases.

\begin{table}[htbp]
\centering
\small
\renewcommand{\arraystretch}{1.3}
\caption{Does stochastic Nesterov accelerate? Summary by noise model}
\label{tab:lines}
\begin{tabularx}{\textwidth}{@{}p{2.3cm}p{2.4cm}Yp{3.2cm}@{}}
\toprule
Noise & Acceleration & Known results & Open \\
\midrule
(N2) persistent & No & The minimax lower bound fixes the order at $\sigma D/\sqrt T$, attained by AC-SA \citep{nemirovski1983,agarwal2012,lan2012}; a.s.\ rates of momentum methods approach $k^{-1/2}$ at best \citep{sebbouh2021,liu2022,attouch2026ravine}; the iterates converge weakly a.s.\ for all $(t_k)$ satisfying $(K_1)$ when $(\gamma_kt_k)\in\ell^2$ and $\gamma_kt_{k-1}^2\to\infty$ \citep{salzo2025iterates}. & Iterate convergence with state-dependent variance \\
\addlinespace
(N5) relative & Yes & Expected $O(1/k^2)$ under suitable noise ratios \citep{vaswani2019fast,gupta2024nesterov,mishkin2024faster}; a.s.\ $o(1/k^2)$ for a two-step-size SNAG under strong growth, for every $\rho\ge1$ \citep{hermant2025gradient}; standard parameters may fail to accelerate \citep{kidambi2018insufficiency,liu2020accelerating}. & Iterate convergence; a.s.\ rates of the single-step-size scheme \\
\addlinespace
(N6) decaying & Yes, if the noise decays fast enough & General-$\alpha_k$ SNAG/SRAG: a.s.\ $O(1/(s_kt_k^2))$, $o$ rates, summable gradients, and iterate convergence \citep{attouch2026ravine,salzo2025iterates}. & $\ell^2$ condition for adaptive $\sigma_k$ \\
\bottomrule
\end{tabularx}
\end{table}

Under persistent noise, acceleration improves only the bias term, so asymptotic acceleration is ruled out. The iterates converge almost surely for every extrapolation sequence satisfying $(K_1)$, with vanishing step sizes and $(\gamma_kt_k)\in\ell^2$ \citep{salzo2025iterates}. Decaying noise is essentially covered by JOTA; for uniformly bounded variance, Salzo's $\ell^2$ condition on $(\gamma_kt_k)$ replaces the $\ell^1$ condition $(K_2^+)$. Under relative noise, acceleration is real: expected $O(1/k^2)$ holds under suitable noise ratios, and a two-step-size SNAG attains a.s.\ $o(1/k^2)$ for every strong growth constant \citep{hermant2025gradient}. Iterate convergence, and the single-step-size scheme with $\alpha_k=1-\alpha/k$, remain open.

% ============================================================
\section{JOTA and the Remaining Questions}
% ============================================================

JOTA studies the regime of general $\alpha_k$ $\times$ decaying noise $\times$ SNAG/SRAG. It extends momentum beyond constant or step-size-driven choices, permits nonincreasing step sizes, including $s_k\to0$, with decaying noise, and transfers the analysis from Nesterov to Ravine through their parameter correspondence.

\begin{table}[htbp]
\centering
\small
\caption{Known results and open gaps for $\alpha_k=1-\alpha/k$}
\label{tab:gaps}
\begin{tabularx}{\textwidth}{@{}p{2.2cm}YYp{1.6cm}@{}}
\toprule
Regime & Known results & Open gap & JOTA \\
\midrule
$\alpha=3$, decaying noise (N6) & a.s.\ $O(1/(s_kk^2))$ under $(K_1)$--$(K_2)$ \citep{attouch2026ravine}; weak a.s.\ iterate convergence \citep{salzo2025iterates} & a.s.\ $o(1/k^2)$ & $O$ rate only \\
$\alpha>3$, decaying noise (N6) & a.s.\ $o(1/(s_kk^2))$, summable gradients, and iterate convergence under $(K_2^+)$ \citep{attouch2026ravine}; $\ell^2$ condition for bounded variance \citep{salzo2025iterates} & $\ell^2$ condition for adaptive $\sigma_k$ & covered \\
$\alpha\ge3$, persistent noise (N2) & rates arbitrarily close to $k^{-1/2}$ \citep{attouch2026ravine,sebbouh2021,liu2022}; weak a.s.\ iterate convergence with vanishing steps \citep{salzo2025iterates} & iterate convergence with state-dependent variance & rate only \\
$\alpha\ge3$, relative noise (N5) & expected $O(1/k^2)$ under suitable ratios \citep{vaswani2019fast,gupta2024nesterov}; a.s.\ $o(1/k^2)$ for the two-step-size SNAG \citep{hermant2025gradient} & iterate convergence; a.s.\ rates of the single-step-size scheme & open \\
\bottomrule
\end{tabularx}
\end{table}

Three questions remain. First, under relative noise and $\alpha_k=1-\alpha/k$ with a single step size, do the iterates converge almost surely, and for which noise ratios $\nu_0(\alpha)$ does the a.s.\ $o(1/k^2)$ rate hold? Gupta et al.\ cover $\alpha=3$ with $\nu<1$ in expectation, and Hermant et al.\ cover every strong growth constant with a two-step-size scheme. Second, does Salzo's $\ell^2$ condition extend from uniformly bounded variance to adaptive variances $\sigma_k$, which contain relative noise? Third, determine whether an a.s.\ $o$ rate is possible at the critical value $\alpha=3$ under decaying noise, where $(K_1^+)$ fails.

% ============================================================
\appendix
\section{Derivation of the One-step Inequality}\label{app:sgd}
% ============================================================

Expanding the squared distance gives
\[
\|x_{t+1}-x^\star\|^2=\|x_t-x^\star\|^2-2s_t\langle\nabla f(x_t)+e_t,x_t-x^\star\rangle+s_t^2\|\nabla f(x_t)+e_t\|^2.
\]
By (N1), the conditional expectation of the cross term with $e_t$ is zero. Convexity gives $\langle\nabla f(x_t),x_t-x^\star\rangle\ge\Delta_t$, and smooth convexity gives $\|\nabla f(x_t)\|^2\le2L\Delta_t$ \citep[Theorem~2.1.5]{nesterov2004}. Hence
\[
\E_t\|x_{t+1}-x^\star\|^2\le\|x_t-x^\star\|^2-s_t(2-2Ls_t)\Delta_t+s_t^2\sigma^2.
\]
When $s_t\le1/(2L)$, this implies \eqref{eq:onestep}.

\bibliographystyle{plainnat}
\bibliography{refs}

\end{document}
